%%%%%%%%%%%%%%%%%%%%%%%%%%%%%%%%%%%%%%%%%%%%%%%%%%%%%%%%%%%%%%%%%%%%%%%%%%%%%%
%
% 01_title.tex
%
%%%%%%%%%%%%%%%%%%%%%%%%%%%%%%%%%%%%%%%%%%%%%%%%%%%%%%%%%%%%%%%%%%%%%%%%%%%%%%

\begin{center}

%%%%%%%%%%%%%%%%%%%%%%%%%%%%%%%%%%%%%%%%%%%%%%%%%%%%%%%%%%%%%%%%%%%%%%%%%%%%%%
%% TITLE
%%%%%%%%%%%%%%%%%%%%%%%%%%%%%%%%%%%%%%%%%%%%%%%%%%%%%%%%%%%%%%%%%%%%%%%%%%%%%%

{\fontsize{62}{68}\selectfont
\bfseries
\color{MOSGold}
THE MATHEMATICS OF THE
MATHEMATICAL OPERATING SYSTEM
}

\vspace{6mm}

{\fontsize{26}{30}\selectfont
\color{MOSGold}
A Cellular--Sheaf Architecture for Structured Machine Cognition
}

\vspace{8mm}

{\Large

\textbf{CHARBEL "THE WIZARD" DARGHAM}

\vspace{2mm}

Lebanese University -- Fanar

\vspace{1mm}

\texttt{cdargham@hotmail.com}

}

\vspace{8mm}

%%%%%%%%%%%%%%%%%%%%%%%%%%%%%%%%%%%%%%%%%%%%%%%%%%%%%%%%%%%%%%%%%%%%%%%%%%%%%%
%% HERO EQUATION
%%%%%%%%%%%%%%%%%%%%%%%%%%%%%%%%%%%%%%%%%%%%%%%%%%%%%%%%%%%%%%%%%%%%%%%%%%%%%%

\begin{eqbox}

\[
\Huge
\boxed{
M=(C,\mathcal{F},s)
}
\]

\[
\Large
\begin{aligned}
C
&=
\text{Stratified Simplicial Complex}
\\[4mm]
\mathcal{F}
&=
\text{Cellular Sheaf}
\\[4mm]
s
&=
\text{Distinguished Global Section}
\end{aligned}
\]

\end{eqbox}

\begin{center}
\begin{tikzpicture}[
  scale=1.2,
  node/.style={circle, fill=MOSRed, draw=MOSGold, line width=0.5pt, minimum size=8pt, inner sep=0pt},
  edge/.style={draw=MOSRed, line width=2pt, opacity=0.8},
  face/.style={fill=MOSRed, opacity=0.3}
]
\coordinate (n1) at (-3,-1);
\coordinate (n2) at (-4,1);
\coordinate (n3) at (-2,3);
\coordinate (n4) at (0,3.5);
\coordinate (n5) at (2,3);
\coordinate (n6) at (4,1);
\coordinate (n7) at (3,-1);
\coordinate (n8) at (1,-2);
\coordinate (n9) at (-1,-2);
\coordinate (n10) at (0,1);
\coordinate (n11) at (-2,1);
\coordinate (n12) at (2,1);
\coordinate (n13) at (0,2.5);

\fill[face] (n1)--(n2)--(n11)--cycle;
\fill[face] (n2)--(n3)--(n11)--cycle;
\fill[face] (n3)--(n4)--(n13)--cycle;
\fill[face] (n4)--(n5)--(n13)--cycle;
\fill[face] (n5)--(n6)--(n12)--cycle;
\fill[face] (n6)--(n7)--(n12)--cycle;
\fill[face] (n7)--(n8)--(n12)--cycle;
\fill[face] (n8)--(n9)--(n10)--cycle;
\fill[face] (n1)--(n9)--(n11)--cycle;
\fill[face] (n11)--(n3)--(n13)--cycle;
\fill[face] (n13)--(n5)--(n12)--cycle;
\fill[face] (n10)--(n12)--(n13)--cycle;
\fill[face] (n10)--(n11)--(n13)--cycle;
\fill[face] (n9)--(n10)--(n11)--cycle;
\fill[face] (n8)--(n10)--(n12)--cycle;

\draw[edge] (n1)--(n2)--(n3)--(n4)--(n5)--(n6)--(n7)--(n8)--(n9)--cycle;
\draw[edge] (n1)--(n11) (n2)--(n11) (n3)--(n11) (n9)--(n11);
\draw[edge] (n3)--(n13) (n4)--(n13) (n5)--(n13);
\draw[edge] (n5)--(n12) (n6)--(n12) (n7)--(n12) (n8)--(n12);
\draw[edge] (n11)--(n13) (n13)--(n12) (n12)--(n10) (n10)--(n11);
\draw[edge] (n9)--(n10) (n8)--(n10);

\foreach \i in {1,...,13} {
  \node[node] at (n\i) {};
}
\end{tikzpicture}
\end{center}

\vspace{6mm}

%%%%%%%%%%%%%%%%%%%%%%%%%%%%%%%%%%%%%%%%%%%%%%%%%%%%%%%%%%%%%%%%%%%%%%%%%%%%%%
%% RESEARCH VISION
%%%%%%%%%%%%%%%%%%%%%%%%%%%%%%%%%%%%%%%%%%%%%%%%%%%%%%%%%%%%%%%%%%%%%%%%%%%%%%

\begin{mosbox}{Research Vision}

\BodyFont

The Mathematical Operating System (MOS) proposes that memory should not
be represented as an unordered vector database but as a structured
mathematical object

\[
M=(C,\mathcal{F},s).
\]

Concepts occupy cells of a simplicial complex,
local information is attached through a cellular sheaf,
and globally coherent cognition is identified with compatible global
sections.

This perspective replaces purely statistical memory with a geometric,
topological and algebraic representation in which reasoning,
verification, semantic organization and structural growth become
intrinsic mathematical operations rather than external heuristics.

\vspace{4mm}

\textbf{Central Philosophy}

\[
\boxed{
\text{Knowledge is Geometry.}
}
\]

\[
\boxed{
\text{Reasoning is Dynamics.}
}
\]

\[
\boxed{
\text{Consistency is Cohomology.}
}
\]

\end{mosbox}

\vspace{8mm}

%%%%%%%%%%%%%%%%%%%%%%%%%%%%%%%%%%%%%%%%%%%%%%%%%%%%%%%%%%%%%%%%%%%%%%%%%%%%%%
%% ARCHITECTURE DIAGRAM
%%%%%%%%%%%%%%%%%%%%%%%%%%%%%%%%%%%%%%%%%%%%%%%%%%%%%%%%%%%%%%%%%%%%%%%%%%%%%%

\begin{tikzpicture}[
scale=1,
every node/.style={
font=\Large,
align=center
}
]

\node[concept] (C)
{
Complex
};

\node[concept,right=4cm of C]
(S)
{
Cellular\\
Sheaf
};

\node[concept,right=4cm of S]
(G)
{
Global\\
Section
};

\node[concept,right=4cm of G]
(H)
{
Coherence
};

\node[concept,right=4cm of H]
(D)
{
Dynamics
};

\draw[arrow] (C)--(S);

\draw[arrow] (S)--(G);

\draw[arrow] (G)--(H);

\draw[arrow] (H)--(D);

%%%%%%%%%%%%%%%%%%%%%%%%%%%%%%%%%%%%%%%%%%%%%%%%%%%%

\node[below=8mm of C]
{
Higher-order Relations
};

\node[below=8mm of S]
{
Local Information
};

\node[below=8mm of G]
{
Compatible Beliefs
};

\node[below=8mm of H]
{
Verification
};

\node[below=8mm of D]
{
Adaptive Memory
};

\end{tikzpicture}

\vspace{10mm}

%%%%%%%%%%%%%%%%%%%%%%%%%%%%%%%%%%%%%%%%%%%%%%%%%%%%%%%%%%%%%%%%%%%%%%%%%%%%%%
%% KEY CONTRIBUTIONS
%%%%%%%%%%%%%%%%%%%%%%%%%%%%%%%%%%%%%%%%%%%%%%%%%%%%%%%%%%%%%%%%%%%%%%%%%%%%%%

\begin{tcolorbox}[%
  colback=MOSDarkRed,
  colframe=MOSGold,
  title=\Large\bfseries Key Contributions,
  coltitle=MOSGold
]

\BodyFont

\begin{itemize}

\item Memory represented as a cellular sheaf over a simplicial complex.

\item Global consistency measured through sheaf Laplacians and
cohomological invariants.

\item Precision-weighted semantic fusion using probabilistic geometry.

\item Topological diagnostics for contradiction localization.

\item Dynamic growth through structural refinement and coarse graining.

\item A mathematical framework intended to augment, rather than replace,
existing reasoning systems.

\end{itemize}

\end{tcolorbox}

\vspace{8mm}

%%%%%%%%%%%%%%%%%%%%%%%%%%%%%%%%%%%%%%%%%%%%%%%%%%%%%%%%%%%%%%%%%%%%%%%%%%%%%%
%% HONEST DISCLAIMER
%%%%%%%%%%%%%%%%%%%%%%%%%%%%%%%%%%%%%%%%%%%%%%%%%%%%%%%%%%%%%%%%%%%%%%%%%%%%%%

\begin{remarkbox}

\BodyFont

\textbf{Current Status}

MOS is an early-stage mathematical research program.

The present work introduces the foundational mathematical architecture,
formal definitions and design philosophy.
Many theoretical components remain under active development and no claim
is made that the system is complete.

The purpose of presenting this work is to invite constructive criticism,
suggestions and collaboration from the mathematical community.

\end{remarkbox}

\end{center}

%%%%%%%%%%%%%%%%%%%%%%%%%%%%%%%%%%%%%%%%%%%%%%%%%%%%%%%%%%%%%%%%%%%%%%%%%%%%%%